% GENERATED FILE. Edit canonical lessons, then run npm run generate:books.
% canonical-source-hash: fnv1a64:9c2136b99796f6a5
% canonical-lessons: MW-C56-ginno, MW-C56-pakarno, MW-C56-katno, MW-C56-rokno, MW-C56-batano

\chapter{To count, To catch, To cut, To stop, To tell}
\label{ch:mw-to-count-to-catch-to-cut-to-stop-to-tell}
\hlchaptermodality{56}

\begin{chapteropening}
\textbf{By the end of this chapter:} \emph{I can say to count, to catch, to cut, to stop and to tell.}

Complete the last lesson of chapter 56: I can say to count, to catch, to cut, to stop and to tell.
\end{chapteropening}

\section[\texorpdfstring{\mw{गिणणो} (giṇṇo)}{गिणणो (giṇṇo)}]{\mw{गिणणो} (giṇṇo) — to count}

\label{lesson:MW-C56-ginno}



\begin{quote}
Before the new one: say the Marwadi for to bring, then the Marwadi for to send.
\end{quote}

\subsection*{You'll want to know: \mw{गिणणो}}
\textbf{\mw{गिणणो}} — \emph{giṇṇo} — \textquotedblleft{}to count\textquotedblright{}.

\begin{cousinweb}
\textbf{\mw{गिणणो}} (\emph{giṇṇo}) is a tadbhava, a word inherited from Sanskrit and reshaped by sound change. It comes from \emph{gaṇayati}, \textquotedblleft{}counts\textquotedblright{}. CDIAL lists Old Marwari \emph{giṇaï} under it (CDIAL 3993). CDIAL thinks the Sanskrit verb itself probably goes back to \emph{gṛṇāti}, \textquotedblleft{}addresses, praises\textquotedblright{}.
\end{cousinweb}

\begin{grammarlens}[title={What you've built}]
One more word for this chapter.
\end{grammarlens}

\subsection*{Your turn}
\begin{itemize}
\raggedright
  \item \emph{Say it:} \emph{giṇṇo}
  \item \emph{Say it:} \emph{giṇṇo}, once more
  \item \emph{Say it:} say \emph{giṇṇo}
  \item \emph{Recall:} say the Marwadi for to bring, then the Marwadi for to send, then say \emph{giṇṇo} again
\end{itemize}

\begin{tcolorbox}[breakable,colback=teal!4,colframe=teal!35!black,title={Before you move on}]
What does \mw{गिणणो} mean? (\textquotedblleft{}To count\textquotedblright{}.) Say it once more.
\end{tcolorbox}
\section[\texorpdfstring{\mw{पकड़णो} (pakaṛṇo)}{पकड़णो (pakaṛṇo)}]{\mw{पकड़णो} (pakaṛṇo) — to catch}

\label{lesson:MW-C56-pakarno}



\begin{quote}
Before the new one: say the Marwadi for to send, then the Marwadi for to count.
\end{quote}

\subsection*{You'll want to know: \mw{पकड़णो}}
\textbf{\mw{पकड़णो}} — \emph{pakaṛṇo} — \textquotedblleft{}to catch\textquotedblright{}.

\begin{cousinweb}
\textbf{\mw{पकड़णो}} (\emph{pakaṛṇo}) is a tadbhava, a word inherited and reshaped by sound change. It comes from \emph{*pakkaḍ}, \textquotedblleft{}seize\textquotedblright{}, a form CDIAL reconstructs from its descendants. CDIAL lists Marwari \emph{pakaṛṇo} under it (CDIAL 7619). CDIAL calls a link with Sanskrit \emph{pragraha}, \textquotedblleft{}holding\textquotedblright{}, obscure and doubtful.
\end{cousinweb}

\begin{grammarlens}[title={What you've built}]
One more word for this chapter.
\end{grammarlens}

\subsection*{Your turn}
\begin{itemize}
\raggedright
  \item \emph{Say it:} \emph{pakaṛṇo}
  \item \emph{Say it:} \emph{pakaṛṇo}, once more
  \item \emph{Say it:} say \emph{pakaṛṇo}
  \item \emph{Recall:} say the Marwadi for to send, then the Marwadi for to count, then say \emph{pakaṛṇo} again
\end{itemize}

\begin{tcolorbox}[breakable,colback=teal!4,colframe=teal!35!black,title={Before you move on}]
What does \mw{पकड़णो} mean? (\textquotedblleft{}To catch\textquotedblright{}.) Say it once more.
\end{tcolorbox}
\section[\texorpdfstring{\mw{काटणो} (kāṭṇo)}{काटणो (kāṭṇo)}]{\mw{काटणो} (kāṭṇo) — to cut}

\label{lesson:MW-C56-katno}



\begin{quote}
Before the new one: say the Marwadi for to count, then the Marwadi for to catch.
\end{quote}

\subsection*{You'll want to know: \mw{काटणो}}
\textbf{\mw{काटणो}} — \emph{kāṭṇo} — \textquotedblleft{}to cut\textquotedblright{}.

\begin{cousinweb}
\textbf{\mw{काटणो}} (\emph{kāṭṇo}) is a tadbhava, a word inherited from Sanskrit and reshaped by sound change. It comes from \emph{kartati}, \textquotedblleft{}cuts\textquotedblright{}. CDIAL lists Old Marwari \emph{kaṭaï} under it (CDIAL 2854). The Old Marwari row means \textquotedblleft{}to be cut\textquotedblright{}; Hindi \emph{kāṭnā} and Gujarati \emph{kāṭvũ}, \textquotedblleft{}to cut\textquotedblright{}, match \textbf{\mw{काटणो}} (\emph{kāṭṇo}) more closely.
\end{cousinweb}

\begin{grammarlens}[title={What you've built}]
One more word for this chapter.
\end{grammarlens}

\subsection*{Your turn}
\begin{itemize}
\raggedright
  \item \emph{Say it:} \emph{kāṭṇo}
  \item \emph{Say it:} \emph{kāṭṇo}, once more
  \item \emph{Say it:} say \emph{kāṭṇo}
  \item \emph{Recall:} say the Marwadi for to count, then the Marwadi for to catch, then say \emph{kāṭṇo} again
\end{itemize}

\begin{tcolorbox}[breakable,colback=teal!4,colframe=teal!35!black,title={Before you move on}]
What does \mw{काटणो} mean? (\textquotedblleft{}To cut\textquotedblright{}.) Say it once more.
\end{tcolorbox}
\section[\texorpdfstring{\mw{रोकणो} (rokṇo)}{रोकणो (rokṇo)}]{\mw{रोकणो} (rokṇo) — to stop}

\label{lesson:MW-C56-rokno}



\begin{quote}
Before the new one: say the Marwadi for to catch, then the Marwadi for to cut.
\end{quote}

\subsection*{You'll want to know: \mw{रोकणो}}
\textbf{\mw{रोकणो}} — \emph{rokṇo} — \textquotedblleft{}to stop\textquotedblright{}.

\begin{cousinweb}
\textbf{\mw{रोकणो}} (\emph{rokṇo}) is a tadbhava, a word inherited and reshaped by sound change. It comes from \emph{*rokk}, \textquotedblleft{}stop\textquotedblright{}, a form CDIAL reconstructs from its descendants. CDIAL gives no Marwari form in this entry (CDIAL 10827); the match is through Hindi \emph{roknā} and Gujarati \emph{rokvũ}. CDIAL says its link with Sanskrit \emph{rudh}, \textquotedblleft{}obstruct\textquotedblright{}, is not clear.
\end{cousinweb}

\begin{grammarlens}[title={What you've built}]
One more word for this chapter.
\end{grammarlens}

\subsection*{Your turn}
\begin{itemize}
\raggedright
  \item \emph{Say it:} \emph{rokṇo}
  \item \emph{Say it:} \emph{rokṇo}, once more
  \item \emph{Say it:} say \emph{rokṇo}
  \item \emph{Recall:} say the Marwadi for to catch, then the Marwadi for to cut, then say \emph{rokṇo} again
\end{itemize}

\begin{tcolorbox}[breakable,colback=teal!4,colframe=teal!35!black,title={Before you move on}]
What does \mw{रोकणो} mean? (\textquotedblleft{}To stop\textquotedblright{}.) Say it once more.
\end{tcolorbox}
\section[\texorpdfstring{\mw{बताणो} (batāṇo)}{बताणो (batāṇo)}]{\mw{बताणो} (batāṇo) — to tell}

\label{lesson:MW-C56-batano}



\begin{quote}
Before the new one: say the Marwadi for to cut, then the Marwadi for to stop.
\end{quote}

\subsection*{You'll want to know: \mw{बताणो}}
\textbf{\mw{बताणो}} — \emph{batāṇo} — \textquotedblleft{}to tell\textquotedblright{}.

\begin{cousinweb}
The origin of \textbf{\mw{बताणो}} (\emph{batāṇo}) is not settled. CDIAL, searched by form and by meaning, has no entry for it or for Hindi \emph{batānā}, \textquotedblleft{}to tell\textquotedblright{}.
\end{cousinweb}

\begin{grammarlens}[title={What you've built}]
One more word for this chapter.
\end{grammarlens}

\subsection*{Your turn}
\begin{itemize}
\raggedright
  \item \emph{Say it:} \emph{batāṇo}
  \item \emph{Say it:} \emph{batāṇo}, once more
  \item \emph{Say it:} say \emph{batāṇo}
  \item \emph{Recall:} say the Marwadi for to cut, then the Marwadi for to stop, then say \emph{batāṇo} again
\end{itemize}

\begin{tcolorbox}[breakable,colback=teal!4,colframe=teal!35!black,title={Before you move on}]
What does \mw{बताणो} mean? (\textquotedblleft{}To tell\textquotedblright{}.) Say it once more.
\end{tcolorbox}
